% BHU PYQ -- History: History of India - Ancient (2025-26 NEP)
\documentclass[11pt,a4paper]{article}
\usepackage[a4paper,top=2.5cm,bottom=2.5cm,left=2.8cm,right=2.8cm,headheight=14pt]{geometry}
\usepackage{amsmath,amssymb}
\usepackage{fontspec}
\setmainfont{Kohinoor Devanagari}
\usepackage{microtype,fancyhdr,enumitem,hyperref,lastpage,parskip}

\hypersetup{colorlinks=true,urlcolor=black,linkcolor=black,pdftitle={HISMJ11: History of India - Ancient -- BHU NEP 2025-26}}

\pagestyle{fancy}\fancyhf{}
\renewcommand{\headrulewidth}{0.4pt}\renewcommand{\footrulewidth}{0.4pt}
\fancyhead[L]{\small \textit{Banaras Hindu University}}
\fancyhead[R]{\small \textit{HISMJ11: History of India - Ancient (2025-26)}}
\fancyfoot[C]{\small Page \thepage\ of \pageref{LastPage}}

\newlist{parts}{enumerate}{2}
\setlist[parts,1]{label=(\alph*),leftmargin=*,itemsep=4pt,topsep=3pt}
\setlist[parts,2]{label=(\roman*),leftmargin=*,itemsep=2pt,topsep=2pt}

\newcommand{\pts}[1]{\hfill \textit{[#1]}}

\begin{document}

\begin{center}
  {\large\bfseries BANARAS HINDU UNIVERSITY}\\[2pt]
  {\normalsize B. A. (Hons.) Semester I Examination, 2025-26 (NEP)}\\[6pt]
  \rule{0.8\linewidth}{0.5pt}\\[6pt]
  {\large\bfseries HISTORY}\\[4pt]
  {\large\bfseries Paper Code: HISMJ11 : History of India -- Ancient}\\[4pt]
  {\normalsize Time: 2:30 Hours \hfill Full Marks: 70}
\end{center}
\rule{\linewidth}{0.4pt}
\smallskip

\noindent \textit{Instructions: This paper comprises three Sections A, B and C. All questions are compulsory. Write your Roll No.\ at the top immediately on receipt of this question paper.}

\smallskip
\noindent \textit{इस प्रश्न-पत्र में तीन खण्ड अ, ब तथा स हैं। सभी प्रश्न अनिवार्य हैं।}

\bigskip

\section*{SECTION -- A / खण्ड -- अ}
\noindent \textit{Note: Very Short Answer Type Questions. Answer each question briefly.}\pts{5 $\times$ 2 = 10}

\noindent \textit{अति लघु उत्तरीय प्रश्न। प्रत्येक प्रश्न का उत्तर दीजिए।}

\begin{enumerate}[label=\textbf{\arabic*.},leftmargin=*,itemsep=8pt]
  \item Explain briefly about the following:\pts{5 $\times$ 2 = 10}
  \begin{parts}
    \item Arthashastra or Harisena / अर्थशास्त्र अथवा हरिषेण
    \item Historical importance of Vedas / वेदों का ऐतिहासिक महत्त्व
    \item Mahapadmananda or Fa-Hien / महापद्मनंद अथवा फाह्यान
    \item Pancha Mahavrata (Five Vows of Jainism) / पञ्च महाव्रत
    \item Kashi Mahajanapada / काशी महाजनपद
  \end{parts}
\end{enumerate}

\bigskip

\section*{SECTION -- B / खण्ड -- ब}
\noindent \textit{Note: Short Answer Type Questions. Answer all questions within 250 words each.}\pts{3 $\times$ 10 = 30}

\noindent \textit{लघु उत्तरीय प्रश्न। सभी प्रश्नों के उत्तर लगभग 250 शब्दों में दीजिए।}

\begin{enumerate}[label=\textbf{\arabic*.},leftmargin=*,start=2,itemsep=12pt]

  \item Throw light on the town-planning, social condition and religious condition of Indus Valley Civilization.\pts{10}
  
  सिन्धु घाटी सभ्यता की नगर योजना, सामाजिक स्थिति तथा धार्मिक स्थिति पर प्रकाश डालिए।

  \begin{center}\textbf{OR / अथवा}\end{center}

  Describe the invasion of Alexander on India. What were its effects?\pts{10}
  
  भारत पर सिकन्दर के आक्रमण का वर्णन कीजिए। इसके क्या प्रभाव हुए?

  \item Throw light on Ashoka's `Dhamma' and describe his religious policy.\pts{10}
  
  अशोक के `धम्म' पर प्रकाश डालिए तथा उसकी धार्मिक नीति का वर्णन कीजिए।

  \begin{center}\textbf{OR / अथवा}\end{center}

  Describe the political achievements of Gautamiputra Satakarni.\pts{10}
  
  गौतमीपुत्र शातकर्णी की राजनीतिक उपलब्धियों का वर्णन कीजिए।

  \item Write a detailed note on the administration of the Mauryan Empire.\pts{10}
  
  मौर्य साम्राज्य के शासन प्रबन्ध (प्रशासन) पर एक विस्तृत टिप्पणी लिखिए।

  \begin{center}\textbf{OR / अथवा}\end{center}

  Discuss the main features of the Vedic society and economy.\pts{10}
  
  वैदिक समाज तथा अर्थव्यवस्था की प्रमुख विशेषताओं की चर्चा कीजिए।

\end{enumerate}

\bigskip

\section*{SECTION -- C / खण्ड -- स}
\noindent \textit{Note: Long Answer Type Questions. Answer all questions in detail.}\pts{2 $\times$ 15 = 30}

\noindent \textit{दीर्घ उत्तरीय प्रश्न। सभी प्रश्नों के विस्तृत उत्तर दीजिए।}

\begin{enumerate}[label=\textbf{\arabic*.},leftmargin=*,start=5,itemsep=14pt]

  \item Write a note on the conquests of Chandragupta Maurya.\pts{15}
  
  चंद्रगुप्त मौर्य की विजयों पर एक टिप्पणी लिखिए।

  \begin{center}\textbf{OR / अथवा}\end{center}

  Throw light on the achievements of Harshavardhana.\pts{15}
  
  हर्षवर्धन की उपलब्धियों पर प्रकाश डालिए।

  \item Who was Pushyamitra Shunga? Describe his political and cultural achievements.\pts{15}
  
  पुष्यमित्र शुंग कौन था? उसकी राजनीतिक एवं सांस्कृतिक उपलब्धियों का वर्णन कीजिए।

  \begin{center}\textbf{OR / अथवा}\end{center}

  Write an essay on any \textbf{one} of the following:\pts{15}
  \begin{parts}
    \item Kalinga King Kharavela and his achievements / कलिंग राजा खारवेल एवं उनकी उपलब्धियाँ
    \item Samudragupta and his conquests / समुद्रगुप्त एवं उनकी विजयें
    \item Cultural development during the Gupta Period / गुप्त काल में सांस्कृतिक विकास
    \item Conquests of Kanishka / कनिष्क की विजयें
  \end{parts}

\end{enumerate}

\end{document}
